\subsection{基环变换}

设 \(K_{0}\) 是含单位元的交换环，\(\rho\) 是将单位元映到单位元的从 \(K_{0}\) 到 K 的同态。设 g 是 K 上的李代数。令 \(g'\) 表示将 g 视为 \(K_{0}\) 上的代数并借助 \(\rho\) 得到的代数（参见第 1 号）。则 \(g'\) 是李代数。g 的子代数（分别为理想）是 \(g'\) 的子代数（分别为理想）。若 a 和 b 是 g 的子模，则括号 {[}a, b{]} 在 g 和 \(g'\) 中相同；因为 {[}a, b{]} 是形如

\(\sum_{i=1}^{n}\left[x_i,y_i\right]\) 的元素组成的集合，其中 \(x_{i}\in a,y_{i}\in b\)。由此 \(\mathcal{D}^p\mathfrak{g} = \mathcal{D}^p\mathfrak{g}'\)、\(\mathcal{C}^p\mathfrak{g} = \mathcal{C}^p\mathfrak{g}'\) 对所有 \(p\) 成立。

子集的中心化子在 \(\mathfrak{g}\) 和 \(\mathfrak{g}'\) 中相同。因此 \(\mathcal{C}_p\mathfrak{g} = \mathcal{C}_p\mathfrak{g}'\) 对所有 \(p\) 成立。

设 \(\mathbf{K}_1\) 是含单位元的交换环，\(\sigma\) 是将单位元映到单位元的从 \(\mathbf{K}\) 到 \(\mathbf{K}_1\) 的同态。设 \(\mathfrak{g}\) 是 \(\mathbf{K}\) 上的李代数。令 \(\mathfrak{g}_{(\mathbf{K}_1)}\) 表示 \(\mathbf{K}_1\) 上通过扩张基环由 \(\mathfrak{g}\) 导出的代数（参见第 1 号）。则 \(\mathfrak{g}_{(\mathbf{K}_1)}\) 是李代数。若 \(\mathfrak{a}\) 是 \(\mathfrak{g}\) 的子代数（分别为理想），则 \(\mathfrak{a}_{(\mathbf{K}_1)}\) 在 \(\mathfrak{g}_{(\mathbf{K}_1)}\) 中的典范像是 \(\mathfrak{g}_{(\mathbf{K}_1)}\) 的子代数（分别为理想）。若 \(\mathfrak{a}\) 和 \(\mathfrak{b}\) 是 \(\mathfrak{g}\) 的子模，则 \(\mathfrak{g}_{(\mathbf{K}_1)}\) 中 \([\mathfrak{a},\mathfrak{b}]_{(\mathbf{K}_1)}\) 的典范像，等于 \(\mathfrak{a}_{(\mathbf{K}_1)}\) 和 \(\mathfrak{b}_{(\mathbf{K}_1)}\) 的典范像的括号。由此，\(\mathcal{D}^p (\mathfrak{g}_{(\mathbf{K}_1)})\) 是 \((\mathcal{D}^p\mathfrak{g})_{(\mathbf{K}_1)}\) 的典范像，而 \(\mathcal{C}^p (\mathfrak{g}_{(\mathbf{K}_1)})\) 是 \((\mathcal{C}^p\mathfrak{g})_{(\mathbf{K}_1)}\) 的典范像。

若 K 是域，\(K_{1}\) 是 K 的扩域，\(\sigma\) 是 K 到 \(K_{1}\) 的典范嵌入，则在通常的同一视下有

\[
\begin{array}{c} [ a, b ] _ {(\mathbf {K} _ {1})} = [ a _ {(\mathbf {K} _ {1})}, b _ {(\mathbf {K} _ {1})} ], \qquad \mathcal {D} ^ {p} (\mathfrak {g} _ {(\mathbf {K} _ {1})}) = (\mathcal {D} ^ {p} \mathfrak {g}) _ {(\mathbf {K} _ {1})}, \\ \mathcal {C} ^ {p} (\mathfrak {g} _ {(\mathbf {K} _ {1})}) = (\mathcal {C} ^ {p} \mathfrak {g}) _ {(\mathbf {K} _ {1})}. \end{array}
\]

这些结果将在第 2 节第 9 号中补充完整。

若 M 是域 K 上的有限维向量空间，则 \(M_{(\mathbf{K}_{1})}\) 是 \(K_{1}\) 上的有限维向量空间，结合代数 \(\mathcal{L}(M_{(\mathbf{K}_{1})})\) 自然地与结合代数 \(\mathcal{L}(M)_{(\mathbf{K}_{1})}\) 相同。因此李代数 \(\mathfrak{gl}(M_{(\mathbf{K}_{1})})\) 自然地与李代数 \(\mathfrak{gl}(M)_{(\mathbf{K}_{1})}\) 相同。

